\subsection{The full metabelian fibres on the rank-two elliptic open}

\begin{statement}{A.171}{Proved}
{Every rational compact fibre, and the exact quadratic lifting locus}
Retain
\[
 \begin{gathered}
 E:\ y^2=x^3-1146^2x,\qquad B=E[2](\Q),\qquad V=E\setminus B,\\
 C=36M,\qquad M=325596479663,\qquad
 \Pi_V=\pi_1(V_{\overline\Q}).
 \end{gathered}
\tag{A.171.1}
\]
Let \(\mathcal S^{\rm met}(V)\) denote geometric conjugacy
classes of continuous sections of \(E_V/\Pi_V''\) locally matching
actual point or cusp sections.  Compactification and the origin of
\(E\) give the actual class map
\[
           \eta:\mathcal S^{\rm met}(V)
                      \longrightarrow H^1_{\rm cts}(\Q,TE).
\tag{A.171.2}
\]
Every rational compact fibre is determined exactly:
\[
 \begin{array}{c|c}
 R\in E(\Q)&\eta^{-1}(\kappa(R))\\\hline
 R\in V(\Q)&\{\text{the metabelian image of the point section at }R\}\\
 R\in B&\{\text{the metabelian images of the full cusp packet at }R\}.
 \end{array}
\tag{A.171.3}
\]
At a boundary point the parametrization of this last fibre by
\(\widehat{\Q^\times}\), using the coordinate \(y\) at finite
two-torsion points and \(u=-x/y\) at the origin, is injective.

The subdomain defined by
\[
                s^*\kappa_2(y/C)=0\quad\text{in }H^1(\Q,\mu_2)
\tag{A.171.4}
\]
consists exactly of the cusp images in (A.171.3) whose parameters
satisfy
\[
                  [\tau_y]_2=[C]_2=[M]_2,
                   \qquad [\tau_u]_2=[-C]_2=[-M]_2,
\tag{A.171.5}
\]
respectively.  Its compact image is precisely
\(\kappa(E[2](\Q))\).  Equivalently the subset of the full
profinite completion of the rank-two group \(E(\Q)\) whose
Kummer classes occur on this locally matching metabelian domain is
exactly the embedded group \(E[2](\Q)\).
For an ordinary rational point the character is the literal
\([y(R)/C]_2\), and it is nonzero.

The first line of (A.171.3) fixes every further completed coordinate
of the section over an ordinary rational compact point.  No
classification of nonrational completed compact classes is asserted.

For the tripod \(U=\mathbf P^1\setminus\{0,1,\infty\}\),
every locally matching section of \(E_U/\Pi''\) with full
abelian pair \((\kappa(r),\kappa(1-r))\),
\(r\in U(\Q)\), is geometrically conjugate to the metabelian
image of the point section at \(r\).  Consequently the full
local point data of \(r\) admit exactly this one global
metabelian section class.
\end{statement}

\begin{proof}
\emph{The compact and affine abelian maps of this actual open curve.}
The map \(V\hookrightarrow E\) induces an arithmetic quotient
\(E_V/\Pi_V''\to\pi_1(E)\).  The origin identifies the latter
with \(TE\rtimes G_\Q\); pulling back the finite multiplication
covers gives (A.171.2).  Choose the actual ordinary rational point
\(R_0=(2028,75348)\) from A.143 as base point for the generalized
Jacobian \(G=J_{E,B}\).  It has the exact sequence
\[
 1\longrightarrow (\mathbf G_m^B)/\mathbf G_m
        \longrightarrow G\longrightarrow E\longrightarrow0,
          \qquad j:V\longrightarrow G,\quad Q\longmapsto[Q-R_0].
\tag{A.171.6}
\]
Here divisor classes are taken with modulus the reduced boundary
\(B\).  This specifies the actual Abel map and all its finite
multiplication pullbacks.  The induced geometric map is
\[
                  \Pi_V^{\rm ab}\xrightarrow{\ \sim\ }TG,
\tag{A.171.7}
\]
as in the primary Proposition A.8(ii),(iii) on the semi-abelian
Albanese of the complement of a normal-crossing divisor
\cite[Proposition A.8(ii),(iii), pp.~232--234]{BIR_MochizukiAlbanese2012}.
Its hypotheses hold literally: \(E\) is smooth projective,
the four boundary components are rational points, and \(R_0\in V(\Q)\).
The source abelianization is already torsion-free: the punctured
genus-one curve has free profinite fundamental group of rank five.
Thus its torsion-free comparison in the cited proposition is the
full comparison (A.171.7).

The compactification sequence is consequently
\[
 0\longrightarrow M_B\longrightarrow\Pi_V^{\rm ab}
          \longrightarrow TE\longrightarrow0,
        \qquad M_B=\widehat\Z(1)[B]/\widehat\Z(1)\sum_{T\in B}e_T.
\tag{A.171.8}
\]
It can also be read directly from the genus-one punctured-surface
presentation: the four boundary meridians have their sole sum
relation, and the remaining two homology generators give the
compact elliptic Tate module.  This describes the exact geometric
maps.  All boundary points are rational, so the arithmetic action
on \(M_B\) is cyclotomic in each of its three free coordinates.

Let \(s\) have compact class \(\kappa(R)\), for
\(R\in E(\Q)\).  Its finite local compact centre is exactly
\(R\).  Indeed compactifying an actual local point or cusp gives
an actual point of \(E(\Q_p)\), and full local Kummer
injectivity makes it equal to \(R\).  This is the formal-subgroup
argument of A.160, and applies at every finite prime.

If \(R\in V\), each actual local section is the point section
at \(R\).  Compare the affine abelian class with this global
point section, and lift a proper conjugator through (A.171.8).
Their difference is in \(H^1(\Q,M_B)\).  It is zero locally:
\(TE^{G_{\Q_p}}=0\), because the local elliptic torsion is
finite, makes \(H^1(\Q_p,M_B)\to H^1(\Q_p,\Pi_V^{\rm ab})\)
injective.  The three split Kummer coordinates are therefore
globally zero by A.169.4.  The full affine abelian image of \(s\)
is the image of the actual point section at \(R\).

If \(R\in B\), each local section is a cusp at \(R\).
Compare with a rational tangent section there instead.  The same
injectivity puts the difference locally in
\(\widehat\Z(1)e_R\subset M_B\).  The inertia-line argument
of A.169.9--11, or the two explicit split quotient coordinates of
A.167.12, recovers one unique global
\(\tau\in\widehat{\Q^\times}\).  Thus the full affine
abelian image is that of an actual full cusp section at \(R\),
with this parameter.  Fix the resulting actual full reference
section \(\sigma\), a point section or a cusp section according
to these two cases.  Lift the affine conjugator, so the original
\(s\) and \(\sigma\) have literally equal abelian images.

\emph{The actual five-dimensional abelian-cover tower.}
For \(m=2j!\), take only the characteristic subgroups
\[
 L_m=\ker(\Pi_V\longrightarrow\Pi_V^{\rm ab}/m\Pi_V^{\rm ab}),
                    \qquad L_m'=\overline{[L_m,L_m]}.
\tag{A.171.9}
\]
These contain \(\Pi_V'\) and satisfy \(\Pi_V''\subset L_m'\).
There is a concrete algebraic model for every arithmetic cover used
here.  Start with the multiplication pullback
\[
       F_m=V\times_{j,G,[m]}G
           =\{(Q,Z):[m]Z=j(Q)\},
                          \qquad F_m\longrightarrow V.
\tag{A.171.10}
\]
It is finite etale of degree \(m^5\), and geometrically connected
by (A.171.7).  The reference section supplies a compatible
continuous cocycle \(c_m:G_\Q\to G[m]\), from its actual affine
abelian class in the splitting based at \(R_0\).
The needed cover \(Y_m\) is the twist of (A.171.10) with the
literal descent action
\[
       \gamma*(Q,Z)=(\gamma Q,\gamma Z-c_m(\gamma)).
\tag{A.171.11}
\]
The cocycle identity verifies this descent, and \([m]c_m=0\)
preserves its defining equation.  Its selected sheet is fixed by
\(\sigma\) and \(s\); its arithmetic subgroup is the full
preimage of their common section in \(E_V/L_m\).
For \(m\mid r\) the actual transition is
\[
                  (Q,Z_r)\longmapsto(Q,[r/m]Z_r).
\tag{A.171.12}
\]
It commutes with (A.171.11), since
\(c_m=[r/m]c_r\).  Thus all finite geometric and arithmetic
cover maps are fixed, rather than just their degrees or isomorphism
classes.

The proper normalization \(C_m\) of \(E\) in this cover has
\(4m^4\) geometric boundary points.  Each of the four primitive
meridians in (A.171.8) has order \(m\) in the deck group, so
the proper ramification index at each boundary is \(m\).
Riemann--Hurwitz gives
\[
                  g(C_m)=1+2m^4(m-1)>1.
\tag{A.171.13}
\]
The fixed sheet of the full reference section \(\sigma\)
supplies a rational point of \(C_m\): it is ordinary above
\(R\) in the point case, and a rational boundary point with its
lifted tangent in the cusp case.  The latter follows by the actual
finite formal-cover fixed-sheet construction used in A.169.
Use it as a proper Abel--Jacobi origin.

The section's image on the proper geometric abelianization of
\(Y_m\) gives a full inverse-Selmer class for
\(\operatorname{Jac}(C_m)\).  Its actual local centres lie in
the finite fibre
\(Z_m=(C_m\to E)^{-1}(R)\), at every finite place.  The
finite fixed-sheet lift gives these local maps, including the real
point or boundary component.  Stoll's Theorem 3.11 applies to the
Abel--Jacobi image of \(Z_m\) without requiring its points to be
rational \cite[Theorem 3.11]{LegStoll2007}.  It gives one
\(Q_m\in Z_m(\Q)\) whose Kummer class is the proper section
class; local full Kummer injectivity and point-injectivity of the
Abel--Jacobi map make all finite local centres equal to this
\(Q_m\).

If \(R\in V\), then \(Q_m\) is ordinary.  Compare with its
actual point section in the affine abelianization.  The remaining
boundary-module class is locally zero.  Choose a finite field
splitting this boundary.  There its module is a sum of copies of
\(\widehat\Z(1)\), so A.169.4 makes the restricted class zero.
The restriction map is injective, by the explicit argument of
A.169.8, because this split module has zero Tate invariants.
The original affine class is therefore the point section at
\(Q_m\).  This argument does not require a rational point in the
boundary of \(Y_m\).

If \(R\in B\), then \(Q_m\) is in that boundary.  The
local classes are in its inertia line.  The nonsplit boundary
argument A.169.9--11 recovers one global tangent parameter and makes
the affine class the image of an actual full cusp section there.
Push this section through the actual cover to \(V\).  The
result is a cusp at \(R\).  Its affine image equals that of the
fixed \(\sigma\); the inertia injection in (A.171.8), together
with zero compact Tate invariants, makes the tangent parameters
equal.  It is therefore the same full cusp class.  In the ordinary
case pushing the point section at \(Q_m\) gives exactly the point
section at \(R\).  In both cases \(s\) and \(\sigma\) are
geometrically conjugate in \(E_V/L_m'\) for every \(m\).

The conjugator sets are nonempty closed subsets of the fixed compact
group \(\Pi_V/\Pi_V''\), nested for divisibility of \(m\).
Their intersection is nonempty.  The exact equality
\(\bigcap_mL_m'=\Pi_V''\) has the same direct proof as
A.169.18: a finite character of \(\Pi_V'\) extends to
\(\Pi_V'W\) for an open normal \(W\) killing its intersection,
and the open image in \(\widehat\Z^5\) contains a sufficiently
divisible \(m\widehat\Z^5\).  The proof uses this actual tower
and does not depend on the rank being two.  The common conjugator
therefore gives equality in \(E_V/\Pi_V''\) itself.
This proves (A.171.3).  At a boundary point the parameter injection
already holds in the affine abelianization, so also in the full
metabelian extension.  Different rational compact points have
different full Kummer classes by local injectivity.

\emph{The quadratic lifting locus and its exact map to the cover.}
The actual cover is A.161's connected finite etale torsor
\[
 D^\circ\longrightarrow V,
       \qquad x=1146(2T-1),\quad y=CW^2,
                   \qquad W^2=y/C.
\tag{A.171.14}
\]
Its geometric subgroup \(H_D\subset\Pi_V\) has index two and
contains \(\Pi_V'\), so
\(\Pi_V''\subset\overline{[H_D,H_D]}\).
A zero class (A.171.4) is precisely a fixed sheet, and gives a
section of the actual subgroup \(E_D/\Pi_V''\).
Its local fixed sheets lift actual points or cusps of \(V\) to
actual points or cusps of \(D^\circ\), by the torsor equation
and its finite formal boundary maps.  It therefore maps to a
locally matching section of
\(E_D/\overline{[H_D,H_D]}\).
A.167 applies through the actual map
\(D^\circ\to U\), \(t=4T(1-T)\): its geometric subgroup is
the same \(\pi_1(D^\circ_{\overline\Q})=H_D\), so the
closed commutator quotient here is exactly the affine cover
quotient used there.  It classifies this section as the image of
one full cusp section of \(D^\circ\).
Under (A.171.14), compactification
sends its class to one of the four actual points
\[
       P_0\mapsto(-1146,0),\quad P_{1/2}\mapsto(0,0),
             \quad P_1\mapsto(1146,0),\quad P_\infty\mapsto O.
\tag{A.171.15}
\]
Thus \(\eta(s)\) is the Kummer class of one of these points.
The already proved boundary fibre in (A.171.3) makes \(s\)
itself cuspidal in the full metabelian quotient.

There is also an explicit map from the tower (A.171.10)--(A.171.12)
to this quadratic cover on this subdomain.  Let
\(\varphi_y:G\to\mathbf G_m\) be the actual homomorphism
defined on degree-zero divisors disjoint from \(B\) by
\[
       \varphi_y\left(\sum_Qn_Q[Q]\right)=\prod_Q y(Q)^{n_Q},
           \qquad \varphi_y(j(Q))=\frac{y(Q)}{y(R_0)}.
\tag{A.171.16}
\]
It descends through the modulus relations: if a rational function
is one at all boundary points, Weil reciprocity pairs its divisor
with \(y\) to give its value on
\(\operatorname{div}(y)=(-1146,0)+(0,0)+(1146,0)-3O\),
which is one.  This specifies the algebraic character and its
normalization.  For even \(m\), choose
\(d\in\overline\Q\) with
\(d^2=y(R_0)/C=2093/M\).  On the geometric multiplication
pullback put
\[
                   W=d\,\varphi_y(Z)^{m/2},
                       \qquad T=\frac{x(Q)/1146+1}{2}.
\tag{A.171.17}
\]
Its square is \(y(Q)/C\), so it gives (A.171.14).
The zero character condition says exactly
\(\varphi_y(c_m(\gamma))^{m/2}=\gamma(d)/d\) in \(\mu_2\).
Under the twisted descent (A.171.11), the factor
\(\varphi_y(c_m)^{-m/2}\) therefore cancels the conjugation
of \(d\); (A.171.17) descends to a \(\Q\)-morphism.
It is the subgroup inclusion \(L_m\subset H_D\), and has
degree \(m^5/2\).  The multiplication transition (A.171.12)
preserves its displayed \(W\).  Thus this tower-to-cover comparison
retains both the arithmetic quadratic twist and its literal root.

Finally the full tangent maps of (A.171.14) are
\[
             \tau_y=C\tau_W^2,\qquad
             \tau_u=-\frac{q_0}{2}\tau_{W/T}^2,
               \qquad q_0=\frac{M}{2\cdot191^2}.
\tag{A.171.18}
\]
The first is the exact function identity.  For the second use
\(u=-(q_0/2)(W/T)^2T/(T-1)\); its last factor has constant
term one and every formal root, as in A.161.11--12.
Since \(C/(q_0/2)=(4\cdot573)^2\), these give precisely
(A.171.5).  Conversely every completed class in (A.171.5) has
the corresponding completed square root: the kernel of reduction
modulo two in \(\widehat{\Q^\times}\) is the image of
squaring, with compatible divided prime valuations as proved in
A.161.  It supplies a full cusp lift.  All four compact torsion
classes occur.  The finite Kummer injections and the continuous
comparison identify the asserted subset of
\(\widehat{E(\Q)}\) exactly, retaining the section domain.
An ordinary point with \([y(R)/C]_2=0\) would give an ordinary
rational point on \(D^\circ\) by (A.171.14), contrary to A.160;
this proves the final nonzero-character assertion.

\emph{The tripod's exact rational-frame fibre.}
Let its full abelian pair be \((\kappa(r),\kappa(1-r))\).
The local cusp equations at zero, one, and infinity are, respectively,
\[
                B=1,\qquad A=1,\qquad B=\kappa(-1)A.
\tag{A.171.19}
\]
Full local Kummer injectivity would give \(r=0\), \(r=1\),
or \(1-r=-r\); all are impossible.  Every finite local centre
is therefore ordinary, and its first coordinate makes it exactly
\(r\).  If this local point data is given initially instead of
the global pair, A.169.4 applied to both coordinates gives precisely
the stated global pair.
Lift the corresponding abelian conjugator through
\(\Pi/\Pi''\to\Pi^{\rm ab}\) and conjugate \(s\), so
its full abelian cocycle is literally that of the point section
\(\sigma_r\), not only the same cohomology class.

For \(n=3j!\), use the actual covers
\[
 \begin{gathered}
 Y_n:\quad rU^n+(1-r)V^n=1,\quad UV\ne0,\qquad t=rU^n,\\
 (U_m,V_m)\longmapsto(U_m^{m/n},V_m^{m/n})\quad(n\mid m),\\
 H_n=\ker(\Pi\longrightarrow\Pi^{\rm ab}/n\Pi^{\rm ab}),
                      \qquad H_n'=\overline{[H_n,H_n]}.
 \end{gathered}
\tag{A.171.20}
\]
They are geometrically connected finite etale covers of degree
\(n^2\).  Their arithmetic subgroups are the preimages of the
common finite section of \(s\) and the point section at \(r\).
Choose the compatible fixed sheets \((U,V)=(1,1)\) of that
point section.  The displayed equations and transition maps then
retain those subgroups and sheets at every level.  Their smooth
proper models are the plane Fermat twists of genus
\((n-1)(n-2)/2\geq1\).  Use the rational ordinary point
\((1,1)\) as the Abel--Jacobi origin.

Since \(\Pi''\subset H_n'\), the metabelian section gives
an actual affine abelian section on \(Y_n\).  Its proper
inverse-Selmer class has all finite local centres in the finite
ordinary fibre over \(r\).  Stoll's Theorem 3.11 gives one
rational point \(Q_n\) of that fibre with this proper Kummer
class.  Local Kummer injectivity and the Abel--Jacobi embedding
make every selected finite local centre equal to \(Q_n\).
Comparison with its point section leaves a class in the boundary
module of \(Y_n\) which is zero at every finite place.  Over a
finite Galois extension splitting that boundary, A.169.4 kills
each split Kummer coordinate.  The zero Tate invariants make
restriction injective, by A.169.8.  Thus this affine class is the
actual point-section image at \(Q_n\), whose pushforward is the
point section at \(r\).

For the fixed point section \(\sigma_r\), form
\[
 \mathcal C_n=\{g\in\Pi/\Pi'':
       g s(\gamma)g^{-1}=\sigma_r(\gamma)
       \text{ in }E_U/H_n'\text{ for all }\gamma\in G_\Q\}.
\tag{A.171.21}
\]
These are nonempty closed subsets of one compact group, nested
along the displayed actual transitions.  The finite-character
argument of A.169.18 gives \(\bigcap_jH_{3j!}'=\Pi''\).
A common conjugator therefore proves equality in the full
metabelian quotient itself.  Conversely the point section is
locally matching and has exactly the given pair, proving the
asserted singleton and the rational-frame conclusion.

The literal quadratic maps, five-dimensional inertia characters,
multiplication-cover degrees and ramification, and twist-cancellation
identity are checked in
\path{code/harmonic_elliptic_metabelian_certificate.py}, with retained
output.  The full metabelian conclusion uses the displayed actual
subgroup tower and the compact conjugator proof.
\end{proof}
